\documentclass[12pt,a4paper,openany]{report}
\usepackage[a4paper,margin=1in]{geometry}
\usepackage[T1]{fontenc}
\usepackage[utf8]{inputenc}
\usepackage{mathptmx}
\usepackage{setspace}
\usepackage{graphicx}
\usepackage{float}
\usepackage{booktabs}
\usepackage{array}
\usepackage{tabularx}
\usepackage{longtable}
\usepackage{xcolor}
\usepackage{hyperref}
\usepackage{caption}
\usepackage{tikz}
\usetikzlibrary{arrows.meta,positioning,shapes.geometric}
\usepackage{listings}
\usepackage{fancyhdr}
\usepackage{titlesec}
\usepackage{enumitem}
\usepackage{url}
\usepackage{placeins}

\hypersetup{colorlinks=true, linkcolor=black, urlcolor=blue, citecolor=black}
\onehalfspacing
\raggedbottom
\setlength{\parindent}{0pt}
\setlength{\parskip}{6pt}
\renewcommand{\arraystretch}{1.15}
\captionsetup{font=bf, labelfont=bf}
\urlstyle{same}

\definecolor{codegray}{RGB}{246,248,250}
\definecolor{codeborder}{RGB}{210,215,220}
\lstdefinestyle{reportcode}{
  backgroundcolor=\color{codegray},
  basicstyle=\ttfamily\small,
  breaklines=true,
  frame=single,
  framerule=0.4pt,
  rulecolor=\color{codeborder},
  columns=fullflexible,
  keepspaces=true,
  showstringspaces=false,
  tabsize=2
}
\lstset{style=reportcode}

\pagestyle{fancy}
\fancyhf{}
\fancyhead[L]{23AIML049}
\fancyhead[R]{Agentic Requirement Analyzer -- Internship Report}
\fancyfoot[L]{CSPIT, CHARUSAT}
\fancyfoot[C]{\thepage}
\fancyfoot[R]{Department of AIML}
\renewcommand{\headrulewidth}{0.4pt}
\renewcommand{\footrulewidth}{0.4pt}
\setlength{\headheight}{15pt}

\titleformat{\chapter}[display]
  {\normalfont\bfseries\centering\Large}
  {CHAPTER \thechapter}
  {0.5em}
  {\Large}
\titlespacing*{\chapter}{0pt}{-5pt}{14pt}
\titleformat{\section}{\normalfont\bfseries\large}{\thesection}{1em}{}
\titleformat{\subsection}{\normalfont\bfseries}{\thesubsection}{1em}{}

\newcommand{\reporttitle}{AGENTIC REQUIREMENT ANALYZER}
\newcommand{\subtitle}{Autonomous Multi-Agent AI System for Software Requirement Analysis and User Story Generation}
\newcommand{\student}{HARI PATEL}
\newcommand{\rollno}{23AIML049}
\newcommand{\internalguide}{Prof. Hardik Jaiswal}
\newcommand{\externalguide}{Mr. Pallav Mamtora}
\newcommand{\externaltitle}{CEO}
\newcommand{\company}{Techmicra IT Solutions}
\newcommand{\hod}{Dr. Nirav Bhatt}


\newcommand{\pre}{\vspace{43pt}}
\newcommand{\post}{\vspace{18.4pt}}
\begin{document}
\pagestyle{fancy}

% ===== original PDF page 17 (report page 8) =====
\setcounter{page}{8}
\setcounter{chapter}{4}\setcounter{section}{3}\setcounter{table}{0}\setcounter{figure}{0}
\section{Risk and Ambiguity Detection Agent}
\vspace{13.4pt}
This agent works differently from the others. A small seed corpus of
known-vague phrases (for example ``quickly'', ``user-friendly'', ``as
needed'') is embedded using a Sentence-Transformers model at module load
time [7] and indexed with FAISS [8]. For every requirement, the agent
embeds the requirement text and searches the index for the closest
known-vague phrases. If the similarity is above a threshold of 0.40, the
requirement and the matched phrase are passed to the LLM together, which
confirms the issue and proposes a concrete rewording -- for example,
replacing ``the system should respond quickly'' with a specific
response-time target. Separately, a rule-based detector classifies each
requirement into security, compliance, integration, performance or data
categories, with a severity and a mitigation, and populates the risks list.

\vspace{8.6pt}
\begin{lstlisting}
query_vector = normalize(embed(requirement))
scores, indexes = faiss_index.search(query_vector, k=3)
for score, index in zip(scores, indexes):
    if score >= similarity_threshold:
        matched_seed = seed_corpus[index]
        ambiguities.append(confirm_with_llm(requirement, matched_seed))
\end{lstlisting}

\vspace{39.4pt}
\section{Technical Code Review Findings}
\vspace{13.4pt}
A detailed review of the implementation and supplied screenshots
identified four engineering defects, all of which were fixed before
submission. First, the ambiguity detector used \texttt{faiss.IndexFlatIP}
without L2-normalized vectors, so scores were raw inner products; seed and
query embeddings are now normalized and scores are true cosine
similarities. Second, every top-$k$ match was forwarded to the LLM; matches
below the 0.40 threshold are now rejected. Third, the \texttt{risks} list
was never populated; risk detection is now separate from ambiguity
detection. Fourth, the acceptance prompt returned values that already
began with ``Given'', ``When'' and ``Then'', while the React renderer
prepended the labels again; the prompt now requests label-free values, the
backend strips any leading label, and the renderer adds each label once.

The upload endpoint was also hardened: only \texttt{.txt}, \texttt{.pdf}
and \texttt{.docx} files are accepted, files are stored under
server-generated names, uploads are limited to 5~MB, and \texttt{/analyze}
accepts only paths inside the upload directory.

% ===== original PDF page 18 (report page 9) =====
\clearpage
\setcounter{section}{5}
\enlargethispage{6pt}
\section{Report Assembly}
\post
The final node in the graph combines the outputs of all six agents into
one structured report containing an executive summary, functional and
non-functional requirements, user stories, acceptance criteria, risks,
ambiguous requirements, and suggestions for improvement. Each completed
run is stored in SQLite so it can be revisited from the frontend's history
view.

\vspace{44pt}
\section{Frontend}
\post
The React and Tailwind frontend provides a drag-and-drop upload form, a
pipeline step indicator, a history sidebar, and a results view with summary
counters, labelled Given/When/Then blocks, ambiguity and risk panels, and
report export. Figure~4.1 shows the upload interface.

\vspace{44pt}
\section{API Endpoints}
\post
The backend exposes the following endpoints for the frontend and manual
testing:

\vspace{27.4pt}
\begin{table}[H]
\centering
\caption{Backend API Endpoints}
\label{tab:api}
\begin{tabularx}{\textwidth}{|p{3.2cm}|X|}
\hline
\textbf{Endpoint} & \textbf{Purpose} \\
\hline
\texttt{GET /health} & Basic health check used to confirm the backend is running. \\
\hline
\texttt{POST /upload} & Accepts a PDF, DOCX or TXT file (max 5~MB) into \texttt{uploads/}. \\
\hline
\texttt{POST /analyze/step1} & Runs parsing, analysis and classification only. \\
\hline
\texttt{POST /analyze} & Runs the full pipeline and saves the report in SQLite. \\
\hline
\texttt{GET /runs} & Returns the saved run history. \\
\hline
\texttt{GET /runs/\{run\_id\}} & Returns one stored analysis result by ID (404 if missing). \\
\hline
\end{tabularx}
\end{table}
\end{document}
